\appendix
\section{LLaVA prompts used}
\label{appendix:llava}

The prompt for zero shot was:

\begin{quote}
You are a housing inspection classifier. Look at the image carefully.

Pick EXACTLY ONE label from: damage, crack, mold, wear, asbestos, no damage

Decision rules (apply in order — stop at first match):
1. asbestos  → flat corrugated or board-like grey/white surface with a matte
               texture; may show thin/thick line markings on the surface.
               NOT rust, NOT peeling paint.
2. crack     → a continuous line-shaped fracture, split, or break in the surface.
               The line must be clearly visible, not just a discoloration.
3. mold      → dark organic spots, green/black patches, mildew, or damp staining
               with fuzzy or spreading edges.
4. damage    → a hole, missing chunk, broken material, dent, or collapse.
               NOT surface aging — the structure itself is broken.
5. wear      → peeling paint, rust, discoloration, stains, or surface aging
               where the structure is still intact.
6. no damage → none of the above are visible.

Tie break rules:
Corrugated grey surface with lines → asbestos (not wear, not no damage)
Thin surface line → crack (not damage)
Brown/orange patches on intact surface → wear (not damage)
Dark fuzzy patches → mold (not wear)

Return ONLY this JSON, nothing else before or after it:
json
{"label": "<one of the 6 labels>", "reason": "<max 10 words>"}
\end{quote}

The prompt of the LLaVA rule based decision prompt was:
\begin{quote}
You are a housing inspection assistant. 

Classify the image into exactly one of these labels:
damage, crack, mold, wear, asbestos, no damage.

Definitions:
- crack: visible cracks, fractures, splits, or line shaped breaks.
- mold: visible mold, mould, mildew, fungal growth, or damp related black spots.
- asbestos: asbestos like material marks, suspicious fibrous material, or asbestos-related surface patterns.
- wear: surface degradation, peeling paint, stains, discoloration, rust, aging, or gradual deterioration.
- damage: other visible damage such as holes, broken surfaces, dents, missing material, fire damage, or water damage.
- no damage: no visible inspection relevant issue.

Rules:
- If a crack is visible, choose crack.
- If mold or mould is visible, choose mold.
- If asbestos like material is visible, choose asbestos.
- Choose no damage when the image only shows a normal wall, floor, ceiling, pipe, surface, or room feature without a clear visible defect. Do not mark minor shadows, texture, lighting, normal edges, or ordinary material patterns as damage.
- Choose no damage only when the area looks fine.

Return only this JSON object:
{
  "label": "one of damage crack mold wear asbestos no damage",
  "summary": "one short sentence"
}
\end{quote}

\section{Use of Artificial Intelligence Tools}
\label{appendix:ai_usage}

During this Thesis, several generative artificial intelligence (AI) tools were used to support software development, learning, experimentation, and thesis improvement. These tools functioned as assistants throughout the research process. All final decisions, implementations, analyses, interpretations, and written content remained the responsibility of the author. I used Vibe Coding.

\subsection{AI Assisted Software Development}

The experimental pipeline was developed using a combination Adjusting code and AI assisted software development techniques. ChatGPT and Claude Code were used throughout the project to support implementation, debugging, and code improvement.

A significant part of the software development process followed a \emph{vibe coding} approach. The author (I) first created a detailed description, conceptual design, and visual representation of the desired functionality. Based on these specifications, AI tools generated an initial implementation. The generated code was subsequently reviewed, modified, tested, debugged, and extended by the author. This process was repeated iteratively until the implementation satisfied the requirements of the research.

Rather than relying on a single AI system, multiple AI models were used to critically evaluate proposed solutions. Code generated by ChatGPT was frequently changed with suggestions from Claude Code, and vice versa. Differences between outputs were analyzed before implementation decisions were made. No AI generated code was incorporated without human review and adaptation.

\subsection{AI as a Learning Tool}

AI tools were also used as educational support throughout the project. When unfamiliar concepts, machine learning techniques, software libraries, or implementation challenges were encountered, ChatGPT was used in a manner similar to a tutor or teaching assistant and functioned as a thesis supervisor.

The tools provided explanations, examples, and alternative approaches that supported the learning process. These explanations were subsequently verified through documentation, scientific literature, experimentation, and practical implementation. AI tools supported understanding but did not replace the author's own learning process, critical thinking, or decision making. It was asked for advise at some point.

\subsection{AI Assisted Thesis Improvement}

Artificial intelligence tools were used to support the organization and improvement of the thesis. Their role was comparable to that of an book editor.

Examples of support activities included:

\begin{itemize}
    \item Structuring reviewer feedback into actionable improvement points
    \item Identifying weaknesses, strengths, and areas requiring clarification
    \item Suggesting improvements to the structure and organization of chapters
    \item Improving readability, grammar, and consistency of existing text
    \item Reviewing draft chapters against the assessment rubric
    \item Highlighting potential gaps in argumentation, methodology, or reporting
\end{itemize}

AI tools were not used to independently generate the scientific contributions, research findings, experimental results, interpretations, or conclusions of this thesis. The author remained responsible for all scientific reasoning, methodological decisions, experimental design choices, data analysis, and interpretation of results.

\subsection{Responsibility and Academic Integrity}

The author maintained full control over all stages of the research process, including dataset selection, model development, implementation, experimentation, evaluation, interpretation, and thesis writing.

Artificial intelligence tools were used as support instruments comparable to coding assistants, educational resources, and editorial feedback mechanisms. The author critically evaluated all generated suggestions and determined whether they were appropriate, correct, and relevant for inclusion in the final work.

Consequently, all scientific claims, research outcomes, interpretations, conclusions, and software implementations presented in this thesis remain the sole responsibility of the author.

%                          
\subsection{ Datasets}
Several publicly available datasets from Roboflow are used for the training process:

\begin{itemize}
    \item Paint damage dataset (67 images)
    \item Crack detection dataset (1551 images)
    \item House damage dataset (3570 images)
    \item Asbestos dataset (1331 images)
    \item Wear dataset (1665 images)
    \item Mold 2 (938 images)
\end{itemize}

These datasets provide additional labeled examples of inspection related defects and help improve model robustness during training.

%\\\\\\\\\\\\\\\\\\\\\\\\\\\\\\\\\\\\\\\\\\\\\\\\\\\\\\\\\\\\\\\\\\\\

\clearpage
\section{Ground Truth Distribution per Property}
\label{appendix:property_ground_truth}

Table~\ref{tab:property_ground_truth} presents the ground truth distribution of the inspection categories for each property before and after the lease period. The overall status indicates whether the property improved, remained unchanged, or deteriorated based on the post lease inspection.
%\\\\\\\\\\\\\\\\\\\\\\\\\\\\\\\\\\\\\\\\\\\\\\\\\\\\
\begin{table*}[ht]
\centering
\caption{Ground truth distribution of the inspection categories for each property.}
\label{tab:property_ground_truth}
\begin{tabular}{llrrrrrrl}
\toprule
Property & Inspection & Damage & Crack & Mold & Wear & Asbestos & No damage & Overall status \\
\midrule

001 & Pre lease  & 3 & 2 & 0 & 3 & 0 & 12 & \multirow{2}{*}{Worse} \\
    & Post lease & 5 & 3 & 0 & 4 & 0 & 8 & \\
\midrule

002 & Pre lease  & 3 & 1 & 0 & 4 & 0 & 12 & \multirow{2}{*}{Worse} \\
    & Post lease & 4 & 2 & 0 & 4 & 0 & 10 & \\
\midrule

003 & Pre lease  & 6 & 4 & 2 & 4 & 1 & 3 & \multirow{2}{*}{Stayed the same} \\
    & Post lease & 6 & 4 & 2 & 4 & 1 & 3 & \\
\midrule

004 & Pre lease  & 7 & 4 & 2 & 3 & 1 & 3 & \multirow{2}{*}{Improved} \\
    & Post lease & 3 & 2 & 1 & 3 & 0 & 11 & \\
\midrule

005 & Pre lease  & 1 & 1 & 0 & 4 & 0 & 14 & \multirow{2}{*}{Improved} \\
    & Post lease & 0 & 0 & 0 & 4 & 0 & 16 & \\
\midrule

006 & Pre lease  & 2 & 0 & 1 & 3 & 0 & 14 & \multirow{2}{*}{Worse} \\
    & Post lease & 4 & 0 & 3 & 4 & 0 & 9 & \\
\midrule

007 & Pre lease  & 3 & 2 & 0 & 3 & 0 & 12 & \multirow{2}{*}{Worse} \\
    & Post lease & 4 & 3 & 0 & 4 & 0 & 9 & \\
\midrule

008 & Pre lease  & 6 & 3 & 1 & 4 & 0 & 6 & \multirow{2}{*}{Improved} \\
    & Post lease & 5 & 3 & 0 & 4 & 0 & 8 & \\

\bottomrule
\end{tabular}
\end{table*}
%\\\\\\\\\\\\\\\\\\\\\\\\\\\\\\\\\\\\\\\\\\\\\\\\\\\\
This table below shows the overview of the property conditions.

\begin{table}[ht]
\centering
\caption{Overview of property conditions.}
\label{tab:property_summary}
\begin{tabular}{lc}
\toprule
Property status & Number of properties \\
\midrule
Improved & 3 \\
Stayed the same & 1 \\
Worse & 4 \\
\bottomrule
\end{tabular}
\end{table}